\reportsection{17}{Icons}{Icon Staging: Extraction and Opacity, Now Shipped Upstream}

\unithead{Mechanism}

\reppar{\code{scripts/staging/icons.sh} (\code{process\_icons}) is the build-time machinery that turns the Windows installer's \code{claude.exe} icon resource into Linux-consumable assets. \code{wrestool -x -t 14 lib/net45/claude.exe -o claude.ico} pulls icon resource type 14 (\code{RT\_GROUP\_ICON}) out of the Windows PE binary; \code{icotool -x claude.ico} expands the multi-resolution container into individual \code{claude\_N\_SIZExSIZEx32.png} files, hard-failing the build if either tool invocation fails. A size-to-index map (\code{16}\,=\,13 through \code{256}\,=\,6) lets deb and rpm packaging install each size into the \code{hicolor} icon theme tree for \code{.desktop} integration; AppImage packaging uses only the 256\texttimes{}256 output, copied to the AppDir's hicolor tree, a top-level PNG, and \code{.DirIcon}. Separately, the function copies the bundled \code{TrayIconTemplate(.png|-Dark.png)} files (and \code{@2x}/\code{@3x} variants) into \code{\$electron\_resources\_dest}, then runs ImageMagick (\code{magick}, falling back to the deprecated \code{convert}) with \code{-channel A -fx 'a>0?1:0' +channel} against each one: any non-zero alpha pixel becomes fully opaque, since the originals are macOS \textasciitilde{}20\%-opacity templates Linux never colorizes. deb, rpm, and AppImage packaging consume the hicolor and tray outputs directly off the filesystem; Nix additionally required a dedicated \code{nix-resources} destination so \code{process\_icons} had somewhere to write, plus a second, un-opacity-processed asar-internal tray-icon copy for the \code{isPackaged=false} code path.}

\unithead{Origin}

\reppar{Pickaxe on \code{wrestool}/\code{icotool} bottoms out at the true origin, \code{d8f4bbe} (2024-12-26, the project's initial commit, as \code{build-deb.sh}): the extraction, size-to-index map, and per-size \code{hicolor} install were present from day one, with no opacity conversion and tray icons copied straight inside the asar. \#134 (\code{47c1889}, merged 2025-11-07, fixing \#122) moved the tray-icon copy out of the asar and onto the filesystem ("Tray icons must be in filesystem (not inside asar) for Electron Tray API to access them"), and created the dedicated Icon Processing section \code{icons.sh} descends from. The opacity conversion itself was born in \code{e5d58f2} (2026-01-19, \#163), a commit shared with \code{tray.sh}: this unit's half added the ImageMagick alpha pass, while the tray section covers the selection-side fix (the light/dark inversion) landed in the same diff.}

\begin{chronology}
2024-12 & \code{d8f4bbe} & Origin: \code{wrestool}/\code{icotool} extraction, size-to-index hicolor map, deb \code{hicolor} install; tray icons copied inside the asar, in \code{build-deb.sh}'s initial commit. \\
2025-11 & \#122, \#134, \code{47c1889} & Tray-icon copy moved from inside the asar to the filesystem \code{\$electron\_resources\_dest}: Electron's \code{Tray} API needs a real filesystem path, not an asar-internal one; dedicated Icon Processing section created. \\
2026-01 & \#163, \code{e5d58f2}; \code{e29635f} & Opacity conversion born (shared commit, tray.sh half is the selection sed), then corrected the next day: the erroneous \code{-negate} removed and the loop widened over \code{TrayIconTemplate*.png} to catch \code{@2x}/\code{@3x} variants. \\
2026-01 & \code{3865848} & Threshold fix: the alpha formula replaced with \code{-fx "a>0?1:0"} after prior attempts left 5\textendash{}20-of-255-alpha edge pixels translucent. \\
2026-01 & \code{29173e9}; \code{86a1822} & Extracted into named \code{process\_icons()} during the 26-function \code{build.sh} refactor; RPM packaging added as a second consumer of the same extraction output. \\
2026-02 & \code{9fe293d}; \code{573f052} & Nix build path wired with its own \code{nix-resources} destination; a second, un-opacity-processed tray-icon copy added inside the asar for Nix's \code{isPackaged=false} code path (@typedrat). \\
2026-03 & \code{0e4a1e7} & First integration-test coverage: \code{tests/test-artifact-\{deb,rpm,appimage\}.sh} check hicolor presence and tray-icon counts (@sabiut). \\
2026-04 & \code{ff4821e} & Moved verbatim into \code{scripts/staging/icons.sh} during the module split. \\
\end{chronology}

\methodbanner{\textbf{Fate under the rebase.} \bdel{} \code{scripts/staging} is deleted wholesale; the official tree ships purpose-made \code{TrayIconLinux(-Dark).png} plus a complete hicolor set at 16\textendash{}256px, which the new packaging copies verbatim.\\[3pt]\textbf{Affects.} All desktops (.desktop hicolor set plus tray icons); the opacity defect specifically hit KDE and Fedora panels (\#163, shared with the tray unit).}
